\section{Scope, source frontier, and principal results}
\label{scope:section}

This article continues the program in the two directories specified
by the project owner: the polylogarithm manuscript and the incoming
research contributions \cite{RepoManuscript,RepoIncoming}. All
repository assertions refer to the pinned revision
\begin{center}\small\ttfamily
c4036d60715916ae72dd18983782fdb182840b86.
\end{center}
The incoming index at that revision contains eleven archives.
We extracted them, reconciled their overlapping research-status
lists, and read the proof dependencies used below. This matters:
an open question in an earlier contribution can already be solved
in a later archive at the same revision.

In particular, the smaller-order-ten and smaller-order-twelve
moment families had already been completed, as had the complete
gap-two and gap-four families, several Stieltjes spectral faces,
and arbitrary mixed even-tail products
\cite{BilinearIncoming,TailKernelsIncoming,OrderedIncoming,LatticeIncoming}.
Those results are inputs or context, and are not claimed anew.
The useful remaining targets were the complete gap-six family,
complex rational harmonic degree, and unequal-rate boundary
collisions. The new theorems address these three targets directly.

\subsection{What changes at this frontier}

\begin{longtable}{@{}>{\raggedright\arraybackslash}p{0.23\textwidth}
 >{\raggedright\arraybackslash}p{0.40\textwidth}
 >{\raggedright\arraybackslash}p{0.29\textwidth}@{}}
\toprule
Question at the pinned revision & Result in this article & Precise scope\\
\midrule
\endfirsthead
\toprule
Question at the pinned revision & Result in this article & Precise scope\\
\midrule
\endhead
Optimal moment prefix for all even pairs with gap six &
Theorem~\ref{gs:main}: the first $q$ centered moments reconstruct
all moments for $p=q-6$, $q\ge8$ even. &
Full formal coefficient rank; an analytic infinite range plus an
exact finite certificate.\\[7pt]
Harmonic degree for complex rational weights &
Theorem~\ref{cr:degree}: degree equals generic pencil rank for
every nonzero complex rational weight. &
Rational dependence on cutoff; only finitely many admissible
cutoff exceptions. No positivity hypothesis.\\[7pt]
Explicit adjacent-pole determinants beyond small rank &
Theorems~\ref{hp:hahn-top} and \ref{hp:all-r-generator}: a
three-by-three Hahn formula and a classical-polylogarithm
generator for every particle count. &
Specified two-pole family, with every endpoint residue included.\\[7pt]
Unequal-rate Stieltjes boundary collisions &
Theorems~\ref{urc:constant-theorem} and \ref{urc:full-theorem}:
all index and argument-derivative orders, with a finite contact
formula for the constant. &
Fixed positive rates, and locally uniform joint spectral
subtraction.\\
\bottomrule
\end{longtable}

Several identity consequences have a particularly short form.
For example, with $H_M=\sum_{k=1}^M1/k$, the positive ensemble
in Section~\ref{hp:section} satisfies
\[\begin{aligned}
 &\sum_{0\le n_1<n_2<n_3<M}
 \frac{\prod_{i<j}(n_j-n_i)^2}{\prod_i(n_i+1)(n_i+2)}\\
 &\qquad=
 -\frac{M(M+1)(M+2)}3H_M^2+
 \frac{4M^2(M+2)}3H_M+
 \frac{M^3(M^2-18M-31)}{12(M+1)}.
\end{aligned}\]
The boundary calculation produces the equally compact ordered
identity
\[
 \operatorname{CT}_{\epsilon\downarrow0}
 \int_0^1\psi'(x+\lambda\epsilon)\psi(x+\mu\epsilon)\,dx
 =\zeta(2)\left(\frac\mu\lambda+\log\frac\mu\lambda\right),
 \qquad\lambda,\mu>0.
\]
Here $\operatorname{CT}$ is explicitly the coefficient of
$\epsilon^0(\log\epsilon)^0$, not the value of a divergent
unsubtracted integral. The article supplies the complete
subtractions as well as this constant.

\subsection{The three mechanisms}

The first mechanism turns a limiting determinant into an exact
all-index result. The incoming Gamma basis has three exponential
scales for gap six. The middle scale loses its leading imaginary
phase, so a one-term Laplace estimate cannot prove the theorem.
Retaining the next term produces a normalized reciprocal-square
Vandermonde. A quantitative error comparison proves nonvanishing
for every $b\ge2048$, and exact recurrence arithmetic covers
$4\le b\le2047$. Estimates are used only to finish an identity and
reconstruction problem; no inequality is presented as the main
research output.

The second mechanism is algebraic. If a symmetric rational matrix
$G(z)$ has a rank-one difference
\[
 G(z)-G(z-1)=w(z)p(z)p(z)^{\mathsf T}
\]
with $w\ne0$ and independent rational components of $p$, then
$G$ cannot be generically singular. Otherwise its rank stabilizes
at large integer cutoffs, equal ranks force its image to remain
fixed, and the update vectors lie in a fixed proper subspace.
That would make a nonzero constant combination of the components
of $p$ vanish identically. This elementary argument works over
every field of characteristic zero, including complex fields with
isotropic vectors.

The third mechanism is local at the singular endpoint. The
Hurwitz recurrence splits each Stieltjes generating function into
one explicit singular power and an analytic function at one.
Only a critical Taylor coefficient of the analytic factor can
contribute to the rate-dependent constant. The singular--singular
term carries a beta or polylogarithmic kernel but contributes
strictly negative powers of the collision parameter. The resulting
constant is therefore much smaller than the full two-shift
integral.

\subsection{Conventions and logical status}

We use $B_1=-1/2$, rising factorials $(a)_k$, and falling
factorials $a^{\underline k}$. Generalized harmonic numbers are
$H_N^{(s)}(a)=\sum_{n=0}^N(n+a)^{-s}$, while the unshifted
notation $H_M$ runs from $1$ to $M$. The polygamma convention is
$\psi^{(r)}=d^{r+1}\log\Gamma/dz^{r+1}$, so
$\psi^{(r)}(z)=(-1)^{r+1}r!\zeta(r+1,z)$ for $r\ge1$.
The generalized Stieltjes constants are fixed by
\[
 \zeta(1+u,z)=\frac1u+
       \sum_{m\ge0}\frac{(-1)^m}{m!}\gamma_m(z)u^m.
\]
Unless otherwise stated, argument derivatives are indicated by
parenthesized superscripts; spectral derivatives are written
explicitly.

Formal harmonic coordinates and formal zeta labels are independent
symbols in their specified coefficient spaces. A rank theorem for
that space is not a theorem of arithmetic independence for its
evaluated constants. Similarly a cutoff generating function is
first an identity on $|x|<1$; a boundary value or an antiderivative
at one is given its own convergence argument.

The article contains four kinds of evidence. Analytic and
algebraic proofs establish the general statements. Exact modular
nonvanishing is part of the gap-six proof on a finite, explicitly
bounded interval. Other exact finite checks verify symbolic
formulas and normalizations, without being substituted for general
proofs. High-precision quadratures are diagnostics, with no claim
that floating-point residuals are certified enclosures.

The methods of Christoffel transformations, Hahn polynomials,
rational discrete residues, Euler sums, and local finite-part
subtraction are established. They are credited to primary
references where used \cite{Krattenthaler2023,DLMF,ChenSinger2012,
Xu2017,XuWang2019,ShpotParis2016}. The claimed advances are the
precise theorems and specializations relative to the inspected
ProveIt frontier. The literature review supports that distinction
but is not an exhaustive historical-priority search. Independent
internal proof audits are included; external peer review and
proof-assistant formalization remain separate tasks.

